\documentclass[11pt,letterpaper]{article}

% ==============================================================================
% PAQUETES FUNDAMENTALES Y CONFIGURACIÓN TIPOGRÁFICA
% ==============================================================================
\usepackage[utf8]{inputenc}
\usepackage[T1]{fontenc}
\usepackage[spanish,es-noshorthands]{babel}
\usepackage[margin=1.75cm, top=1.75cm, bottom=1.75cm, headheight=14pt]{geometry}
\usepackage{amsmath,amssymb,amsfonts,bm}
\usepackage{booktabs}
\usepackage{tabularx}
\usepackage{enumitem}
\usepackage{microtype}
\usepackage{xcolor}
\usepackage{fancyhdr}
\usepackage{lastpage}
\usepackage{float}

% ==============================================================================
% PALETA INSTITUCIONAL UIS
% ==============================================================================
\definecolor{uisdark}{HTML}{1A1D20}
\definecolor{uispurple}{HTML}{3D2B56}
\definecolor{uisblue}{HTML}{0D47A1}
\definecolor{uisborder}{HTML}{CBD5E1}
\definecolor{uisgray}{HTML}{64748B}

% ==============================================================================
% HIPERVÍNCULOS
% ==============================================================================
\usepackage{hyperref}
\hypersetup{
    colorlinks=true,
    linkcolor=uispurple,
    citecolor=uisblue,
    urlcolor=uispurple,
    pdftitle={Programa Sintético: Análisis No Lineal de Estructuras},
    pdfauthor={Prof. Orlando Arroyo, Ph.D. - UIS},
    pdfsubject={Programa Sintético - Posgrados},
    pdfkeywords={Análisis No Lineal, OpenSeesPy, Moodle UIS, Pushover, Método N2}
}

% ==============================================================================
% ENCABEZADOS Y PIES DE PÁGINA
% ==============================================================================
\pagestyle{fancy}
\fancyhf{}
\fancyhead[L]{\small\color{uisgray}\textsf{Análisis No Lineal de Estructuras}}
\fancyhead[R]{\small\color{uisgray}\textsf{Programa Sintético}}
\fancyfoot[L]{\small\color{uisgray}\textsf{Escuela de Ingeniería Civil}}
\fancyfoot[C]{\small\color{uisgray}\textsf{Posgrados}}
\fancyfoot[R]{\small\color{uisgray}\textsf{Pág. \thepage\ de \pageref{LastPage}}}
\renewcommand{\headrulewidth}{0.4pt}
\renewcommand{\footrulewidth}{0.4pt}

% ==============================================================================
% COMIENZO DEL DOCUMENTO
% ==============================================================================
\begin{document}

% ==============================================================================
% PÁGINA 1: IDENTIDAD INSTITUCIONAL, ECOSISTEMA Y OBJETIVOS FORMATIVOS
% ==============================================================================
\enlargethispage{1.5\baselineskip}

\begin{center}
    {\scshape\Large Universidad Industrial de Santander}\\[0.05cm]
    {\large Escuela de Ingeniería Civil}\\[0.04cm]
    {\normalsize\bfseries Posgrados}\\[0.12cm]
    {\color{uispurple}\hrule height 1.8pt}\vspace{0.12cm}
    
    {\huge\bfseries\color{uisdark} Análisis No Lineal de Estructuras}\\[0.08cm]
    {\large\bfseries\color{uispurple} Programa Sintético y Guía de Contenidos}\\[0.12cm]
    
    {\color{uispurple}\hrule height 1.0pt}\vspace{0.14cm}
    
    \small
    \begin{tabularx}{\textwidth}{@{} l X l X @{}}
        \textbf{Profesor:} & Prof. Orlando Arroyo, Ph.D. & \textbf{Periodo:} & Semestre Académico \\
        \textbf{Nivel:} & Posgrados & \textbf{Modalidad:} & Teórico-Práctica Computacional \\
        \textbf{Software:} & OpenSeesPy (\texttt{openseespy}) & \textbf{Plataforma:} & Moodle UIS
    \end{tabularx}
\end{center}

\vspace{0.02cm}

\section{Descripción General y Ecosistema Digital}

La asignatura de \textbf{Análisis No Lineal de Estructuras} aborda los principios mecánicos, formulaciones matemáticas y herramientas computacionales necesarias para modelar y evaluar la respuesta inelástica y la capacidad sismorresistente de estructuras de concreto reforzado ante sismos severos, superando las simplificaciones del análisis elástico tradicional.

El desarrollo del curso se articula mediante el siguiente ecosistema de herramientas digitales oficiales:

\begin{enumerate}[leftmargin=16pt, topsep=1.5pt, itemsep=2pt, parsep=0pt, label=\textbf{\arabic*.}]
    \item \textbf{Plataforma Web Interactiva del Curso (Landing Page):} \newline
    \href{https://odarroyo.github.io/analisis_no_lineal/}{\small\texttt{https://odarroyo.github.io/analisis\_no\_lineal/}} \newline
    {\small Laboratorios numéricos interactivos: simulaciones de dependencia de trayectoria (\textit{path-dependence}), residuo de desequilibrio $\mathbf{P} - \mathbf{F}_{\text{int}}(\mathbf{u}) = \mathbf{0}$, algoritmos de Newton-Raphson e inspección matricial paso a paso.}

    \item \textbf{Portal Complementario de Concreto Reforzado:} \newline
    \href{https://odarroyo.github.io/curso_concreto_reforzado/}{\small\texttt{https://odarroyo.github.io/curso\_concreto\_reforzado/}} \newline
    {\small Recursos técnicos sobre diseño sísmico seccional, análisis momento-curvatura ($M-\phi$), diagramas de interacción $P-M$ y detallamiento dúctil seccional.}

    \item \textbf{Entorno Computacional Oficial --- OpenSeesPy:} \newline
    {\small Modelación numérica en Python mediante la biblioteca \textbf{OpenSeesPy} (\texttt{openseespy}), con soporte de librerías científicas (\texttt{NumPy}, \texttt{SciPy}, \texttt{Matplotlib}, \texttt{opsvis}, \texttt{vfo}) para resolución no lineal y posprocesamiento.}

    \item \textbf{Aula Virtual Institucional (Moodle UIS):} \newline
    {\small Espacio oficial para la gestión académica, foros técnicos y recepción de entregables del proyecto semestral y evaluaciones.}
\end{enumerate}

\vspace{0.05cm}

\section{Objetivos Formativos}

\subsection{Objetivo General}
Capacitar al estudiante de posgrado en los fundamentos mecánicos, algoritmos numéricos de equilibrio y técnicas avanzadas de modelado inelástico de estructuras de concreto reforzado en \textbf{OpenSeesPy}, orientados a la evaluación sismorresistente basada en desempeño.

\subsection{Objetivos Específicos}
\begin{itemize}[leftmargin=16pt, topsep=1pt, itemsep=1.2pt, parsep=0pt]
    \item Comprender los principios matemáticos del equilibrio inelástico, la dependencia de trayectoria y la convergencia de algoritmos iterativos (Newton-Raphson, control de desplazamiento).
    \item Caracterizar y modelar fuentes de no linealidad geométrica ($P-\Delta$, $P-\delta$, corotacional) y material (concreto confinado y acero con efecto Bauschinger).
    \item Formular e implementar modelos para pórticos (plasticidad concentrada y fibras) y muros estructurales (láminas multicapa, macroelementos MVLEM y SFI-MVLEM).
    \item Aplicar el análisis Pushover multimodal y el Método N2 Modificado para determinar el punto de desempeño y verificar estados límite (ASCE/SEI 41 y Eurocódigo 8).
\end{itemize}

\clearpage

% ==============================================================================
% PÁGINA 2: CONTENIDOS TEMÁTICOS (LISTADO SINTÉTICO)
% ==============================================================================

\section{Contenidos Temáticos (Listado Sintético)}

A continuación se presenta el listado sintético de los 7 contenidos estructurados para la asignatura:

\begin{enumerate}[leftmargin=18pt, topsep=2pt, itemsep=3.5pt, parsep=0pt, label=\textbf{\arabic*.}]
    \item \textbf{¿Qué es el análisis no lineal?}
    \begin{itemize}[leftmargin=12pt, topsep=1pt, itemsep=1pt]
        \item Principios del comportamiento inelástico y dependencia de trayectoria (\textit{path-dependence}).
        \item No conmutatividad de combinaciones de carga y pérdida del principio de superposición.
        \item Formulación de la ecuación de equilibrio residual: $\mathbf{r}(\mathbf{u}) = \mathbf{P} - \mathbf{F}_{\text{int}}(\mathbf{u}) = \mathbf{0}$.
        \item Algoritmos iterativos de Newton-Raphson (estándar, modificado y BFGS).
        \item Criterios y tolerancias de convergencia (fuerza, desplazamiento y energía).
        \item Estrategias de control de carga, control de desplazamiento y longitud de arco.
    \end{itemize}

    \item \textbf{Tipos de análisis estructural: Lineal vs. No lineal, Estático vs. Dinámico}
    \begin{itemize}[leftmargin=12pt, topsep=1pt, itemsep=1pt]
        \item Clasificación: Lineal estático (ELF), modal espectral, estático no lineal (pushover) y dinámico no lineal paso a paso (NL-THA).
        \item Hipótesis fundamentales, costos computacionales y límites de aplicabilidad.
        \item Significado físico y limitaciones del factor de reducción de fuerza por ductilidad $R$.
        \item Criterios normativos de selección del método (NSR-10, ASCE 7, Eurocódigo 8).
    \end{itemize}

    \item \textbf{Fuentes de no linealidad}
    \begin{itemize}[leftmargin=12pt, topsep=1pt, itemsep=1pt]
        \item \textbf{3.1 No linealidad de material:} Tracción, fluencia del refuerzo, aplastamiento del concreto, ablandamiento (\textit{softening}), pérdida de adherencia y degradación cíclica.
        \item \textbf{3.2 No linealidad geométrica:} Efectos de segundo orden $P-\Delta$ y $P-\delta$, matriz de rigidez geométrica $\mathbf{K}_G$ y transformaciones geométricas en OpenSeesPy (Lineal, $P-\Delta$, Corotacional).
        \item \textbf{3.3 Modelos no lineales de materiales:} Modelos uniaxiales de concreto confinado y no confinado de Mander (\texttt{Concrete01}, \texttt{Concrete02}) y modelo de Menegotto-Pinto para acero de refuerzo (\texttt{Steel02}).
    \end{itemize}

    \item \textbf{Formulaciones no lineales para elementos de pórticos de concreto reforzado}
    \begin{itemize}[leftmargin=12pt, topsep=1pt, itemsep=1pt]
        \item \textbf{4.1 Modelos de plasticidad concentrada:} Rótulas plásticas de longitud nula ($M-\theta$), longitud de plastificación $L_p$ y leyes histeréticas con deterioro cíclico tipo IMK (\texttt{ModIMKPinching}).
        \item \textbf{4.2 Modelos de fibras:} Discretización seccional en fibras, elementos basados en desplazamientos (\texttt{dispBeamColumn}) vs. basados en fuerzas (\texttt{forceBeamColumn}) con regularización de energía.
    \end{itemize}

    \item \textbf{Formulaciones para elementos de edificios de muros}
    \begin{itemize}[leftmargin=12pt, topsep=1pt, itemsep=1pt]
        \item \textbf{5.1 Formulaciones Shell Multilayer:} Elementos finitos laminares multicapa 2D con capas ortótropas bajo teoría de campo de compresión modificada (MCFT).
        \item \textbf{5.2 Formulación de fibras:} Modelo de columna ancha (\textit{wide column}) y discretización en fibras.
        \item \textbf{5.3 Formulaciones de MVLEM y SFI-MVLEM:} Macroelemento de múltiples líneas verticales (MVLEM) y macroelemento con interacción corte-flexión (SFI-MVLEM) en OpenSeesPy.
    \end{itemize}

    \item \textbf{Análisis pushover (estático no lineal)}
    \begin{itemize}[leftmargin=12pt, topsep=1pt, itemsep=1pt]
        \item Aplicación secuencial de cargas de gravedad y empuje lateral bajo control de desplazamiento.
        \item Patrones de carga lateral (uniforme, triangular, modal) y curva de capacidad ($V_b$ vs. $\Delta_{\text{techo}}$).
        \item Transformación cinemática al oscilador equivalente SDOF ($\Gamma$, $M^*$) e identificación del mecanismo de colapso.
    \end{itemize}

    \item \textbf{Método N2 modificado}
    \begin{itemize}[leftmargin=12pt, topsep=1pt, itemsep=1pt]
        \item Formulación en el plano aceleración-desplazamiento (ADRS: $S_a - S_d$) y bilinealización por energía.
        \item Factor de reducción inelástica $R_\mu$ y cálculo del Punto de Desempeño $d^*$.
        \item Extensión modificada para efectos de modos superiores y torsión acoplada.
        \item Verificación de estados límite de desempeño (IO, LS, CP) según ASCE/SEI 41-17 y Eurocódigo 8.
    \end{itemize}
\end{enumerate}

\clearpage

% ==============================================================================
% PÁGINA 3: METODOLOGÍA, EVALUACIÓN Y BIBLIOGRAFÍA
% ==============================================================================

\section{Metodología de Enseñanza y Esquema de Evaluación}

La metodología combina exposiciones conceptuales de fundamentación teórico-mecánica con sesiones prácticas guiadas de programación y modelación no lineal en \textbf{OpenSeesPy}. La plataforma \textbf{Moodle UIS} centraliza el seguimiento académico, foros y gestión de entregables.

El sistema de evaluación se articula en torno a un \textbf{Proyecto Práctico de Modelación No Lineal} desarrollado en dos entregas consecutivas que concentran el \textbf{90\%} de la calificación final, complementado por un \textbf{Quiz Final} conceptual correspondiente al \textbf{10\%}.

\vspace{0.10cm}

\begin{table}[H]
\centering
\caption{Esquema oficial de evaluación de la asignatura (Posgrados).}
\label{tab:evaluacion_sintetica}
\small
\setlength{\tabcolsep}{4.5pt}
\renewcommand{\arraystretch}{1.15}
\begin{tabularx}{\textwidth}{l X c c}
\toprule
\textbf{Actividad Evaluativa} & \textbf{Alcance, Enfoque y Entregables en Moodle UIS} & \textbf{Porcentaje} & \textbf{Semanas} \\
\midrule
\textbf{Proyecto: Entrega 1} & 
\textbf{Modelación Constitutiva Seccional y de Componentes en OpenSeesPy:} \newline
Calibración de leyes no lineales de concreto (Mander) y acero (Menegotto-Pinto). Modelación comparativa de pórticos (plasticidad concentrada vs. fibras) y muros (MVLEM / SFI-MVLEM) con validación experimental. & 
\textbf{45\%} & Sem. 8--9 \\
\midrule
\textbf{Proyecto: Entrega 2} & 
\textbf{Modelo Estructural Global, Pushover y Método N2 Modificado:} \newline
Ensamblaje del modelo inelástico completo en OpenSeesPy. Análisis pushover multimodal, curva de capacidad, transformación ADRS, punto de desempeño con N2 modificado y evaluación de estados límite (ASCE 41 / Eurocódigo 8). & 
\textbf{45\%} & Sem. 15--16 \\
\midrule
\textbf{Quiz Final} & 
\textbf{Evaluación Conceptual de Fundamentos y Algoritmos:} \newline
Prueba individual sobre principios matemáticos del análisis inelástico, residuo de equilibrio, algoritmos de Newton-Raphson, matrices tangentes, convergencia y modelos de elementos finitos. & 
\textbf{10\%} & Sem. 16 \\
\midrule
\textbf{Total} & \textbf{Evaluación Integral de la Asignatura} & \textbf{100\%} & \\
\bottomrule
\end{tabularx}
\end{table}

\vspace{0.05cm}

\section{Bibliografía Principal}

\begingroup
\footnotesize
\begin{enumerate}[leftmargin=16pt, topsep=1pt, itemsep=0.8pt, parsep=0pt, label={[\arabic*]}]
    \item \textbf{Chopra, A. K. (2020).} \textit{Dynamics of Structures: Theory and Applications to Earthquake Engineering}. 5th Edition, Pearson Education.
    \item \textbf{Priestley, M. J. N., Calvi, G. M., \& Kowalsky, M. J. (2007).} \textit{Displacement-Based Seismic Design of Structures}. IUSS Press, Pavia, Italy.
    \item \textbf{Powell, G. H. (2010).} \textit{Nonlinear Analysis of Structures for Seismic Design}. Computers and Structures Inc., Berkeley, CA.
    \item \textbf{Fajfar, P. (2000).} ``A Nonlinear Analysis Method for Performance-Based Seismic Design''. \textit{Earthquake Spectra}, 16(3), 573--592.
    \item \textbf{Kolozvari, K., Orakcal, K., \& Wallace, J. W. (2015).} ``Modeling of Cyclic Shear-Flexure Interaction in RC Structural Walls. I: Theory''. \textit{Journal of Structural Engineering}, 141(5), 04014135.
    \item \textbf{Mander, J. B., Priestley, M. J. N., \& Park, R. (1988).} ``Theoretical Stress-Strain Model for Confined Concrete''. \textit{Journal of Structural Engineering}, 114(8), 1804--1826.
    \item \textbf{Menegotto, M., \& Pinto, P. E. (1973).} ``Method of analysis for cyclically loaded RC frames including changes in geometry and non-elastic behavior of elements''. \textit{IABSE Symposium}, Lisbon.
    \item \textbf{ASCE/SEI 41-17 (2017).} \textit{Seismic Evaluation and Retrofit of Existing Buildings}. American Society of Civil Engineers, Reston, VA.
    \item \textbf{Eurocode 8 (EN 1998-1) (2004).} \textit{Design of structures for earthquake resistance}. CEN, Bruselas.
    \item \textbf{OpenSees Documentation Team (2024).} \textit{OpenSeesPy Command Manual}. \url{https://openseespydoc.readthedocs.io/}
\end{enumerate}
\endgroup

\vspace{0.20cm}
\begin{center}
    {\color{uispurple}\hrule height 1.2pt}\vspace{0.08cm}
    {\footnotesize\color{uisgray} Universidad Industrial de Santander --- Escuela de Ingeniería Civil --- Posgrados --- Bucaramanga, Colombia}
\end{center}

\end{document}
